\chapter{Evidence and technical notes}
\label{app:notes}

This appendix preserves the detail that an implementer or scholar may need
after reading the main argument. It is not a project diary. Search ledgers,
raw responses, source hashes, and package manifests remain in the repository's
evidence record; this appendix states only the conclusions relevant to the
proposal.

\section{How to read the evidence status}

The literature and software review was assembled from primary papers,
systematic reviews, official documentation, and internal case records. Source
identity and claim support were treated as separate questions. A DOI match
establishes bibliographic identity; it does not establish that the source
supports a sentence. A package's official documentation establishes an
available function; it does not establish effectiveness on technical
documents.

At the current run, ten of fifteen evidence gates pass. The bounded
RePEc/IDEAS specialist slice, live known-item recall checks, and Crossref DOI
identity checks are complete. The outstanding items are independent screening
and adjudication, full-text anchors and design-specific appraisal for
load-bearing studies, held-out package benchmarks, and independent external
review. The proposal is therefore an author-repaired research request, not a
completed systematic review or an efficacy report.

\begin{boundarybox}{Claim classes used in this proposal}
\begin{description}
\item[Source result.] The cited source explicitly reports the finding at the
  inspection depth available.
\item[Local implication.] The product consequence follows from the source
  result and the stated workflow, without adding an efficacy claim.
\item[Project hypothesis.] The proposed workflow may reduce burden or improve
  reconstruction; the feasibility and comparative studies must test it.
\item[Open question.] The current evidence does not decide the issue and names
  the next observation required.
\end{description}
\end{boundarybox}

\section{The compact internal case record}

The internal cases are product requirements, not external efficacy evidence.
They are useful because they contain a failure, a repair, and a reader
consequence rather than a collection of style opinions.

\begin{table}[H]
\centering
\small
\begin{tabularx}{\textwidth}{@{}p{0.21\textwidth}L p{0.28\textwidth}@{}}
\toprule
Case & Observed pattern & Design consequence \\
\midrule
Funding proposal & A compiled and repeatedly reviewed draft left a reader
unfamiliar with the project unable to state purpose or requested action. &
The reading brief and acceptance by an unfamiliar reader belong in the product,
not in an after-the-fact review. \\
Model note repair & A sequence of revisions preserved equations while removing
private labels and making the comparison explicit. & Protected spans and
located findings that the author can accept or reject are feasible targets. \\
Accepted technical units & Small units passed a named reader gate after
substantive explanation, while whole-document acceptance was not claimed. &
Measure the unit of reader decision and do not generalize a local repair. \\
Long monograph review & Mathematical detail survived, but repeated structure
and administrative language increased the reading burden. & Diagnostics must
serve a running argument and cannot be accepted on the basis of surface scores. \\
\bottomrule
\end{tabularx}
\caption{Internal cases motivate the product boundary; they do not estimate its
effect.}
\label{tab:internal-cases}
\end{table}

\section{Protected objects and allowed changes}

The protected-source model distinguishes object preservation from prose
preservation. The following default is appropriate for the first fixtures:

\begin{longtable}{@{}p{0.25\textwidth}p{0.28\textwidth}p{0.38\textwidth}@{}}
\toprule
Object & Default check & Permitted author action \\
\midrule
\endfirsthead
\toprule
Object & Default check & Permitted author action \\
\midrule
\endhead
Equation body & Exact normalized hash and rendered comparison. & Change only
with an explicit explanation and a rechecked result. \\
Numerical literal and unit & Token and unit comparison. & Change with an
explanation of the new estimate or convention. \\
Citation key and bibliography identity & Key resolution and metadata match. &
Replace or remove after human source verification. \\
Label, cross-reference, and table membership & Resolve after each build. &
Move when the rendered relation remains correct. \\
Claim or assumption & Author-declared span and substantive diff. & Rewrite,
narrow, or remove with an author decision. \\
Quotation, code, and data field & Exact or format-specific comparison. &
Change only under the source's disclosure and ownership rules. \\
\bottomrule
\end{longtable}

\section{The first software benchmark}

The benchmark corpus should contain at least one valid and one malformed case
for each of the following: custom LaTeX macros, displayed mathematics,
citations, labels, tables, comments, code blocks, Markdown emphasis, and
source-to-render locations. For each parser candidate, record whether the
protected object is recovered, whether a finding points to the same source
span on replay, and whether a deliberate change is classified correctly.

For diagnostics, human reviewers label a frozen sample by category and burden.
The benchmark reports precision, abstention, runtime, and review minutes. It
does not report a single "human" score. For model-assisted findings, the
benchmark also records order sensitivity, length sensitivity, provider and
version, and the proportion of cases in which the model abstains.

\section{Evidence that would change the recommendation}

The recommendation to fund the MVP should be revised if any of the following
new evidence appears:

\begin{itemize}
\item a direct competitor already provides the complete source-to-reader path
  with lower measured burden;
\item protected comparisons cannot preserve equations, citations, or claims
  across the target formats;
\item independent readers find that the review packet makes documents harder
  to understand or increases their time;
\item the live document cannot be processed within its privacy and disclosure
  boundary; or
\item the comparative design cannot observe a credible reader-acceptance
  endpoint.
\end{itemize}

Conversely, a successful feasibility packet is not sufficient to claim
efficacy. It earns the right to run the comparative study and to price the
larger decision.
